\section{Exact finite graph certificates}\label{sec:certificates}
The preceding lemmas reduce the required physical bounds to finite rational
polynomial arithmetic, but this reduction alone is not a seed. We describe
how the stored certificates are tied to the actual graphs.

On a torus, enumerate sites row-major, with $0$ fixed. The cubic entries in
\eqref{eq:K} have integer coefficient arrays after scaling by $64$.
Expanding \eqref{eq:q2}--\eqref{eq:q6} produces exact rational polynomials.
For the quartic term, every support containing $0$ is present, including those
whose coefficient vanishes. There are respectively
$\binom{15}{3}=455$, $\binom{31}{3}=4495$, and $\binom{63}{3}=39711$ supports.
Each has three matching labels. Signed fiber sums are collected \emph{before}
absolute values or Gram weights are applied. Their degree-at-most-ten
polynomials are represented by eleven integer coefficients with common
normalization $524288$. The three displayed envelope dictionaries have
43, 508, and 462 entries; these are coefficient classes, not counts of physical
supports or numerical samples.

\begin{proposition}[Whole-interval certification]\label{prop:interval}
Every stored coefficient bound and residual norm bound used below holds for
all real $0\le g\le1/10$. The certificate hypotheses are discharged by exact
finite graph identities and rational inequalities.
\end{proposition}
\begin{proof}[Certificate argument]
For each coefficient dictionary entry $c(g)$, the proof supplies a rational
$u$ and checks positivity of $u-c(g)$ and $u+c(g)$ on the entire interval.
After the change $x=10g$, represent a polynomial as
\[
 p(x)=\sum_{k=0}^{d} b_k\binom dk x^k(1-x)^{d-k}.
\]
Nonnegative rational Bernstein coefficients imply $p(x)\ge0$ for $0\le x\le1$.
The basis-conversion identity and coefficient inequalities are proved exactly.
The pair squared-norm polynomials and the collected signed quartic Gram
polynomials are treated in the same way; the latter are not replaced by the
sum of absolute local terms.

Separately, exact table theorems identify each actual graph coefficient with
its dictionary polynomial, up to its stored sign. Row validity, increasing
codes, and exact binomial cardinality prove that the rooted support enumeration
is a bijection. Bounded chunks and balanced concatenation retain every row and
every boundary comparison. They change evaluation cost, not the mathematical
condition. The covariance cache is proved equal to $|K_{rs}(1/10)|+\tau_A$
for every pair by actual displacement translation. The entrywise monotonicity
proved above bounds $L(g)$ by this endpoint covariance on the whole interval. When a translated support
is reordered, the same explicit permutation transports its matching label,
coefficient, and covariance weight. Therefore \eqref{eq:rootcount} applies to
both the Gram and weighted-mean sums without assuming enumeration invariance.
Exact integer aggregates then prove the rational ceilings.

The generated arrays are machine-readable appendices in the formal source.
Ordinary Python generation and diagnostic checks help produce these witnesses;
neither their exit status nor an imported JSON number is an axiom. The accepted
Lean terms prove graph membership, interval inequalities, and aggregate values.
\end{proof}

For clarity, the quantities $F_{A,d}$ in Table~\ref{tab:vacuum} bound
\emph{squared faithful vacuum norms}. They are multiplied by $M_{A,d}$ before
square roots are taken. The actual physical component bounds follow from
\eqref{eq:projection} and \eqref{eq:amplify}. Only areas 16 and 32 need residual
norms; area 64 needs its mean but no Temple variance estimate.
\begin{table}[ht]
\centering\small
\begin{tabular}{c cc}\toprule
 & $A=16$ & $A=32$\\\midrule
$F_{A,2}$ & $\dfrac{2871661602999}{819200000000000000}$ & $\dfrac{10257117031323}{2621440000000000000}$\\[9pt]
$F_{A,4}$ & $\dfrac{10086255523899002668131}{419430400000000000000000000}$ & $\dfrac{61781177216998702643223}{1342177280000000000000000000}$\\[9pt]
$\ell_{A,6}$ & $\dfrac{92232993647601}{655360000000000000}$ & $\dfrac{2956223510791297}{10485760000000000000}$\\[9pt]
$F_{A,6}$ & $27\ell_{16,6}^2$ & $27\ell_{32,6}^2$\\\bottomrule
\end{tabular}
\caption{Exact finite vacuum certificates.}\label{tab:vacuum}
\end{table}

\begin{proposition}[Actual cubic trial bounds]\label{prop:trialbounds}
Let $\mu_A=\ip{\psi_K}{H^{\rm C}\psi_K}/\norm{\psi_K}^2$.
Uniformly on $[0,1/10]$, the actual unnormalized bounds
\[
 |\ip{\psi_K}{(H^{\rm C}-e_A)\psi_K}|\le m_A,\qquad
 \norm{(H^{\rm C}-e_A)\psi_K}\le\varepsilon_A
\]
hold with the rational ceilings in Table~\ref{tab:trial}. Hence
$|\mu_A-e_A|\le m_A$, and the normalized variance is at most
$\varepsilon_A^2$ where a norm bound is supplied.
\end{proposition}
\begin{proof}
Use the exact residual identity, the faithful norm and positive mean estimates,
and the whole-interval graph certificates. For $d=2,4,6$ choose rational
$\eta_{A,d}>0$ with $F_{A,d}M_{A,d}\le\eta_{A,d}^2$; exact squaring checks give
$\sum_d\eta_{16,d}<10982/10^6$ and
$\sum_d\eta_{32,d}<20172/10^6$.
The finite Wick coefficient aggregates multiplied by the separately bounded
positive norm $M_{A,0}$ give $m_A$. Finally $\norm{\psi_K}\ge1$ and centering
at the actual Rayleigh quotient cannot increase the normalized residual
variance. This last step does not identify $e_A$ with $\mu_A$.
\end{proof}
\begin{table}[ht]\centering
\begin{tabular}{c c c c}\toprule
Box & $M_{A,0}$ upper bound & $m_A$ & $\varepsilon_A$\\\midrule
$4\times4$ & $1070873423/10^9$ & $167/10^6$ & $10982/10^6$\\
$8\times4$ & $286527607/250000000$ & $346/10^6$ & $20172/10^6$\\
$8\times8$ & $328014697/250000000$ & $800/10^6$ & not required\\\bottomrule
\end{tabular}
\caption{Uniform actual trial ceilings. Norms are Euclidean Hilbert norms.}\label{tab:trial}
\end{table}
